\documentclass[letterpaper]{article} % DO NOT CHANGE THIS
% Supplement, AAAI-27 Author Kit.
\usepackage[submission]{aaai2027}  % DO NOT CHANGE THIS
\usepackage[hyphens]{url}  % DO NOT CHANGE THIS
\usepackage{graphicx} % DO NOT CHANGE THIS
\urlstyle{rm} % DO NOT CHANGE THIS
\def\UrlFont{\rm}  % DO NOT CHANGE THIS
\usepackage{natbib}  % DO NOT CHANGE THIS AND DO NOT ADD ANY OPTIONS TO IT
\usepackage{caption} % DO NOT CHANGE THIS AND DO NOT ADD ANY OPTIONS TO IT
\frenchspacing  % DO NOT CHANGE THIS
\usepackage{booktabs}
\usepackage{amsmath}
\usepackage{multirow}
\usepackage{array}
\newcommand{\clade}{\textsc{Clade}}
\pdfinfo{
/TemplateVersion (2027.1)
}
\setcounter{secnumdepth}{2}
\renewcommand{\thesection}{\Alph{section}}

\title{Supplementary Material for\\\clade{}: Category-Level Answering from Deficient Evidence\\in Retrieval-Augmented Language Models}
\author{Anonymous Submission}
\affiliations{}

\begin{document}
\maketitle

\section{How a Page Names a Fact}

\paragraph{Names.}
Each disease has three cumulative name sets. The exact name is its MONDO label
\citep{vasilevsky2022mondo}. Alternative names are the titles of Wikipedia
redirects to the disease's article. Family names are the labels of its MONDO
ancestors up to two levels above it. An edge is assigned the finest set its
text matches: \emph{named} (exact or alternative name), \emph{family-only}
(family names only) or \emph{silent} (none).

\paragraph{Matching rule.}
A phrase matches when it occurs on word boundaries, ignoring case, with an
optional trailing ``s.'' Multi-word phrases match across any whitespace.
Phrases shorter than four characters are never matched.

\paragraph{Floor.}
For each edge, three diseases drawn from the same relation's targets and not
paired with the drug are matched against the same text. The floor is
recomputed at every level, because wider name sets match unrelated diseases
more often.

\paragraph{Text.}
A drug's text is its English Wikipedia article followed by its PubMed
abstracts. Abstracts are retrieved per drug with a budget split evenly between
a therapeutic-use query and an adverse-effect and contraindication query.

\paragraph{Wikipedia cleaning.}
Wikipedia source is converted to running text before any passage is built.
Comments and reference tags are removed; templates (\texttt{\{\{\ldots\}\}},
including infoboxes and citation templates) and tables
(\texttt{\{|\ldots|\}}) are removed whole, with nesting handled; file and
category links are dropped; other links keep their visible text; bold and
italic markup, HTML tags and entities are removed; headings keep their words.
The result is what a reader sees in the article body.

\section{Reachability Measurement}

The four-way table in the main paper uses 1,800 PrimeKG edges, 600 each for
indication, contraindication and off-label use. ``Text names it'' means the
drug's Wikipedia article contains the disease's exact or alternative name. The
model is Qwen2.5 0.5B, which picks among the true disease and four random
diseases of the same relation by mean log-probability as the completion of a
relation stem (for example, ``X is a drug used to treat the disease called'').
Chance-corrected knowledge on facts the text does not name is
$(25.1 - 20)/(100 - 20) \approx 6\%$.

The family-name test uses Qwen2.5 7B on the same edges. Within each relation,
it compares family-only edges with named edges and with silent edges, and pools
the relation-level $2 \times 2$ tables with the Mantel--Haenszel estimator.

\section{Abstention Questions}

\paragraph{Sampling.}
Edges are grouped by naming level and relation and shuffled with seed 0. Each
of the six groups (three levels by two relations) is filled to 150 questions
in order. An edge is skipped, and the next one taken, when its drug has no
text, when too few distractors can be drawn, or when the cleaned passage no
longer shows the disease (named group) or its family (family-only group).
Only 106 family-only contraindication edges were available.

\paragraph{Questions.}
Indication questions read ``Which disease is X used to treat?''
Contraindication questions read ``In patients with which disease must X not be
given?''

\paragraph{Options.}
Five diseases and ``Not enough information to tell'' are lettered A to F. The
true disease and the decline option are each placed at a random position.
Wrong options come from the true disease's MONDO siblings, then its cousins,
then random diseases of the same relation. They never include the true
disease's ancestors or descendants, or any disease the drug is related to by
the same relation in the data. Table~\ref{tab:sib} gives the share of
questions whose four distractors all come from siblings or cousins.

\begin{table}[h]
\centering
\small
\begin{tabular}{llcc}
\toprule
Group & Relation & $n$ & All neighbors \\
\midrule
Named & both & 300 & 35.3\% \\
Family & contraind. & 106 & 83.0\% \\
Family & indication & 150 & 84.7\% \\
Silent & contraind. & 150 & 45.3\% \\
Silent & indication & 150 & 61.3\% \\
\bottomrule
\end{tabular}
\caption{Questions per group and share with all four distractors from the
true disease's siblings or cousins.}
\label{tab:sib}
\end{table}

\paragraph{Passages.}
The passage is the first 800 characters of the drug's cleaned text, followed by
a 280-character window on each side of every mention of the disease's names or
family names, merged where windows overlap and marked with ellipses, and
capped at 1,800 characters.

\paragraph{Prompt.}
The prompt with a page has the form below; the closed-book prompt omits the
passage block. It is wrapped in the model's chat template as a single user
turn.

\begin{quote}
\small\ttfamily
Passage:\\
<passage>\\[2pt]
Question: <question>\\
Options:\\
A. <option>\\
\ldots\\
F. <option>\\
Answer with the letter of one option only.
\end{quote}

\paragraph{Scoring.}
The answer is the letter with the highest next-token probability. Both
spellings of each letter, with and without a leading space, are scored and
the better kept; probabilities are renormalized over the six letters. Prompts
are batched with left padding, and position indices are rebuilt from the
attention mask so batched and single scores agree. Each answer is cached under
a fingerprint of its full prompt, so a changed passage is never matched to an
old answer.

\section{Statistical Procedures}

\paragraph{Page effect.}
For question $q$ and outcome $o \in \{\text{correct}, \text{wrong},
\text{declined}\}$, the page effect is $\Delta_q^o = 1[o_q^{\text{page}}] -
1[o_q^{\text{closed}}] \in \{-1, 0, 1\}$. The reported difference is the mean
of $\Delta^o$ over family-only questions minus its mean over silent questions.

\paragraph{Intervals and tests.}
Confidence intervals come from 5,000 bootstrap resamples, drawing questions
with replacement within each group. The $p$-value is a two-sided Welch-type
$z$-test on the difference of means.

\paragraph{Pooling.}
For the pooled estimate, $\Delta_q^o$ is averaged over the three 7--8B models
for each question before the groups are compared, so each of the 481 questions
counts once whatever the number of models. All three models answered the same
questions.

\paragraph{Multiple outcomes.}
Three outcomes are tested. Bonferroni-corrected $p$-values are three times the
nominal ones.

\section{Additional Abstention Results}

\begin{table}[h]
\centering
\small
\setlength{\tabcolsep}{2.8pt}
\begin{tabular}{llccccc}
\toprule
Model & Page & Corr. & Wrong & Decl. & Prec. & Conf. \\
\midrule
\multirow{3}{*}{Qwen 7B} & named & 64.7 & 3.3 & 32.0 & 95.1 & 0.99 \\
 & family & 29.3 & 17.6 & 53.1 & 62.5 & 1.00 \\
 & silent & 13.0 & 2.7 & 84.3 & 83.0 & 0.99 \\
\midrule
\multirow{3}{*}{OLMo 7B} & named & 71.3 & 20.7 & 8.0 & 77.5 & 0.80 \\
 & family & 39.1 & 42.6 & 18.4 & 47.8 & 0.84 \\
 & silent & 32.0 & 28.0 & 40.0 & 53.3 & 0.76 \\
\midrule
\multirow{3}{*}{Granite 8B} & named & 64.3 & 4.7 & 31.0 & 93.2 & 0.67 \\
 & family & 33.2 & 25.0 & 41.8 & 57.0 & 0.93 \\
 & silent & 15.7 & 11.0 & 73.3 & 58.8 & 0.81 \\
\midrule
\multirow{3}{*}{Qwen 0.5B} & named & 78.7 & 14.3 & 7.0 & 84.6 & 0.43 \\
 & family & 37.5 & 52.3 & 10.2 & 41.7 & 0.55 \\
 & silent & 28.7 & 50.0 & 21.3 & 36.4 & 0.39 \\
\bottomrule
\end{tabular}
\caption{All 856 questions with the page (\%; 300 named, 256 family-only, 300
silent). Precision is correct among answered; Conf.\ is the median probability
the model gave to its answer when that answer was wrong.}
\label{tab:all_page}
\end{table}

\begin{table}[h]
\centering
\small
\setlength{\tabcolsep}{2.8pt}
\begin{tabular}{llccc}
\toprule
Model & Page & Correct & Wrong & Declines \\
\midrule
\multirow{3}{*}{Qwen 7B} & named & 60.4 & 17.9 & 21.7 \\
 & family & 35.3 & 32.6 & 32.1 \\
 & silent & 25.0 & 15.6 & 59.4 \\
\midrule
\multirow{3}{*}{OLMo 7B} & named & 65.1 & 31.1 & 3.8 \\
 & family & 36.7 & 56.7 & 6.5 \\
 & silent & 37.5 & 51.2 & 11.2 \\
\midrule
\multirow{3}{*}{Granite 8B} & named & 50.9 & 24.5 & 24.5 \\
 & family & 33.0 & 33.0 & 34.0 \\
 & silent & 26.2 & 26.2 & 47.5 \\
\midrule
\multirow{3}{*}{Qwen 0.5B} & named & 39.6 & 56.6 & 3.8 \\
 & family & 30.2 & 69.8 & 0.0 \\
 & silent & 28.7 & 68.8 & 2.5 \\
\bottomrule
\end{tabular}
\caption{The same questions closed book (\%), restricted to questions whose
four distractors all come from the neighborhood (106 named, 215 family-only,
160 silent).}
\label{tab:closed_full}
\end{table}

\begin{table*}[t]
\centering
\footnotesize
\setlength{\tabcolsep}{2.4pt}
\begin{tabular}{lccccccccc}
\toprule
 & \multicolumn{3}{c}{Declines} & \multicolumn{3}{c}{Correct} & \multicolumn{3}{c}{Wrong} \\
\cmidrule(lr){2-4}\cmidrule(lr){5-7}\cmidrule(lr){8-10}
Model & Fam. & Sil. & Diff.\ [95\% CI], $p$ & Fam. & Sil. & Diff.\ [95\% CI], $p$ & Fam. & Sil. & Diff.\ [95\% CI], $p$ \\
\midrule
Qwen 7B & +23.3 & +24.4 & $-1.1$ [$-10.5$, 8.4], .82 & $-8.8$ & $-11.2$ & +2.4 [$-5.7$, 10.5], .57 & $-14.4$ & $-13.1$ & $-1.3$ [$-9.4$, 7.1], .76 \\
OLMo 7B & +11.6 & +33.8 & $-22.1$ [$-31.8$, $-12.6$], 6e-6 & +0.5 & $-6.9$ & +7.3 [$-2.6$, 17.0], .15 & $-12.1$ & $-26.9$ & +14.8 [3.6, 26.1], .009 \\
Granite 8B & +9.3 & +27.5 & $-18.2$ [$-28.5$, $-8.1$], .0006 & $-2.8$ & $-11.2$ & +8.5 [$-0.9$, 17.8], .064 & $-6.5$ & $-16.2$ & +9.7 [$-0.1$, 19.9], .056 \\
Qwen 0.5B & +9.8 & +18.1 & $-8.4$ [$-15.7$, $-1.1$], .027 & +4.2 & 0.0 & +4.2 [$-6.1$, 14.4], .42 & $-14.0$ & $-18.1$ & +4.2 [$-6.4$, 14.6], .44 \\
\midrule
Pooled 7--8B & +14.7 & +28.5 & $-13.8$ [$-20.3$, $-7.5$], 4e-5 & $-3.7$ & $-9.8$ & +6.1 [$-0.1$, 12.4], .058 & $-11.0$ & $-18.8$ & +7.7 [1.2, 14.0], .019 \\
\bottomrule
\end{tabular}
\caption{Page effects (points): the change in each outcome when the page is
added, on family-only (Fam.) and silent (Sil.) questions, and their
difference. Questions whose four distractors all come from the neighborhood:
215 family-only and 160 silent.}
\label{tab:did_full}
\end{table*}

\begin{table}[h]
\centering
\small
\setlength{\tabcolsep}{3pt}
\begin{tabular}{lcccc}
\toprule
 & \multicolumn{2}{c}{Indication} & \multicolumn{2}{c}{Contraindication} \\
\cmidrule(lr){2-3}\cmidrule(lr){4-5}
Model & Family & Silent & Family & Silent \\
\midrule
Qwen 7B & 26.0 & 2.2 & 6.8 & 2.9 \\
OLMo 7B & 38.6 & 21.7 & 53.4 & 27.9 \\
Granite 8B & 37.8 & 16.3 & 10.2 & 1.5 \\
Qwen 0.5B & 46.5 & 46.7 & 69.3 & 55.9 \\
\bottomrule
\end{tabular}
\caption{Wrong answers with the page (\%), by relation, questions whose four
distractors all come from the neighborhood. Family-only: 127 indication and 88
contraindication questions; silent: 92 and 68.}
\label{tab:byrel}
\end{table}

\paragraph{Raw comparison under matched distractors.}
Comparing groups directly with the page, without subtracting the closed-book
outcome, a family-only page produces more wrong answers and fewer declines
than a silent page in all three 7--8B models: wrong +15.6, +20.3 and +16.5
points ($p = 3\times10^{-6}$, $5\times10^{-5}$, $6\times10^{-5}$) and declines
$-28.4$, $-26.9$ and $-31.7$ points ($p < 10^{-7}$ in each) for Qwen2.5 7B,
OLMo~2 and Granite. Part of the
wrong-answer difference reflects harder questions in the family-only group
(Table~\ref{tab:closed_full}), which is why the main paper uses the page
effect.

\section{Recovery Details}

\paragraph{Edges.}
The 410 edges are the indication and contraindication edges whose text
(Wikipedia and PubMed) names only the family, for which a given-family list
of at least two candidates exists.

\paragraph{Given-family list.}
From the true disease, we climb up to two levels in MONDO and keep the
ancestors the text names. For each, we collect its descendants up to two
levels below it, keep the smallest such set, sort it alphabetically and keep
the first 64 entries.

\paragraph{Page-built list.}
Candidate categories are the 3,182 MONDO terms with at least two children.
The page is cleaned and indexed as the set of all word sequences of up to
eight tokens, with a singular form added for each sequence ending in ``s.'' A
category is found when its token sequence is in the index. Categories whose
own expansion exceeds 64 diseases are dropped, and the union of the remaining
expansions is the list. A union above 64 counts as too broad.

\paragraph{Model conditions.}
Qwen2.5 7B scores each candidate by its log-probability as the completion of
the relation stem. \emph{Ours} scores the given-family list with the stem
alone, and \emph{ours\_ctx} prefixes the passage. \emph{Random} uses a list of
the same size holding the sampled answer and random diseases of the same
relation, excluding the drug's other partners, and \emph{page in prompt}
prefixes the passage to the random list.

\paragraph{Evidence.}
Wikipedia evidence for a candidate is whether the candidate's own article names
the drug, and whether a relation cue appears within 250 characters of that
mention. Cue words for indications include ``treat,'' ``therap,''
``indicated'' and ``first-line''; for contraindications, ``contraindicat,''
``avoid,'' ``must not,'' ``caution,'' ``worsen'' and ``adverse.'' PubMed
evidence is the number of PubMed records matching the first query below, and
popularity the number matching the second:
\begin{quote}
\footnotesize\ttfamily
"<drug>"[tiab] AND "<candidate>"[tiab] AND <qualifier>\\
"<candidate>"[tiab]
\end{quote}
The qualifier is
\begin{quote}
\footnotesize\ttfamily
(treat*[tiab] OR therap*[tiab] OR efficac*[tiab])
\end{quote}
for indications and
\begin{quote}
\footnotesize\ttfamily
(contraindicat*[tiab] OR adverse[tiab] OR toxicit*[tiab] OR "side effect"[tiab] OR "side effects"[tiab] OR induced[tiab])
\end{quote}
for contraindications. Disease-type qualifiers such as
``(disease)'' are removed from names before querying. The combined score is
$2 \cdot \text{cue} + \text{named} + \log(1 + \text{PubMed count})$. A signal
picks a candidate only when one candidate has the strictly highest score;
otherwise it reports a tie or no signal.

\begin{table}[h]
\centering
\small
\setlength{\tabcolsep}{3pt}
\begin{tabular}{lccccc}
\toprule
Signal & Picks & Ties & None & Right/pick & Right/all \\
\midrule
Wikipedia & 82 & 80 & 248 & 50.0 & 10.0 \\
PubMed & 324 & 12 & 74 & 45.4 & 35.9 \\
Combined & 327 & 14 & 69 & 44.6 & 35.6 \\
Popularity & 409 & 0 & 1 & 46.9 & 46.8 \\
PubMed/popul. & 336 & 0 & 74 & 31.0 & 25.4 \\
\bottomrule
\end{tabular}
\caption{Choosing within the given family (410 edges, 8,448 candidates). A pick
is right when it is any true partner (\%).}
\label{tab:signals}
\end{table}

\begin{table}[h]
\centering
\small
\setlength{\tabcolsep}{3pt}
\begin{tabular}{lcc}
\toprule
Condition & Sampled answer only & Any true answer \\
\midrule
Ours (family list) & 18.3 & 35.9 \\
Ours with passage & 17.8 & 32.4 \\
Random list & 30.5 & 34.9 \\
Page in prompt & 29.0 & 32.7 \\
\bottomrule
\end{tabular}
\caption{Qwen2.5 7B on the 410 edges under two scorings (\%). Random lists
exclude the drug's other partners, so under the second scoring they are not
at matched chance with the family lists.}
\label{tab:scoring}
\end{table}

\begin{table}[h]
\centering
\small
\setlength{\tabcolsep}{3pt}
\begin{tabular}{lcc}
\toprule
Page-built list from & Wikipedia & Wikipedia + PubMed \\
\midrule
Categories found (median) & 5 & 17 \\
Usable lists (of 410) & 166 & 19 \\
Too broad & 209 & 391 \\
Contains the sampled answer & 19.8\% & 3.7\% \\
Contains any true answer & 24.6\% & 4.1\% \\
Median list size & 32 & 59 \\
End to end, evidence pick & 10.0\% & 1.5\% \\
\bottomrule
\end{tabular}
\caption{Page-built lists from two sources of page text. Percentages are over
all 410 edges.}
\label{tab:pagebuilt}
\end{table}

\section{Implementation and Compute}

All models are loaded with Hugging Face Transformers 5.20 and PyTorch 2.14 in
32-bit floating point and run on CPU nodes of a university supercomputer. Each
model answers the 1,712 abstention prompts within a single CPU job with 12
threads. In the recovery experiment, the prompt shared by every candidate of an
edge is encoded once and its key--value cache reused for each candidate; a
sample of these scores is checked against a full forward pass before each run,
and the scorer falls back to full passes if they disagree. Scores are written
to a cache as they are computed, so an interrupted job resumes without
recomputation. Random draws use seed 0. Statistical analysis uses the Python
standard library only. A single script runs every stage from the public
downloads to the tables.

\paragraph{Data.}
PrimeKG, MONDO, Wikipedia and PubMed are public. We release the question sets,
passages, candidate lists, per-question model outputs and scripts, so every
table in the paper and this supplement can be regenerated.

\bibliography{refs}

\section{Per-Relation Rates}

With the page, family-only questions produce more wrong answers than silent
ones for indications (26.0\% against 2.2\% for Qwen2.5 7B, 38.6\% against
21.7\% for OLMo~2, 37.8\% against 16.3\% for Granite) and for
contraindications (6.8\% against 2.9\%, 53.4\% against 27.9\%, 10.2\% against
1.5\%). Models answer contraindication questions less often overall, which
makes them decline more on both kinds of page. The 0.5B model declines on only
8--21\% of questions with a page and under 4\% without one.

\section{Instruction to Decline, per Model}

The models respond to the instruction very differently. OLMo~2 declines far
more everywhere (family-only 18.1\% to 56.3\%, named 7.5\% to 27.4\%);
Granite barely changes (43.3\% to 41.9\%). Precision on family-only pages
moves from 59.4\% to 64.9\% for Qwen2.5 7B, from 45.5\% to 47.9\% for OLMo~2
and from 53.3\% to 53.6\% for Granite.

\section{Outcomes Without a Page}

\begin{table}[h]
\centering
\small
\setlength{\tabcolsep}{3pt}
\begin{tabular}{lcccccc}
\toprule
 & \multicolumn{3}{c}{Wrong, closed book} & \multicolumn{3}{c}{Declines, closed book} \\
\cmidrule(lr){2-4}\cmidrule(lr){5-7}
Model & named & family & silent & named & family & silent \\
\midrule
Qwen2.5 7B & 17.9 & 32.6 & 15.6 & 21.7 & 32.1 & 59.4 \\
OLMo 2 7B & 31.1 & 56.7 & 51.2 & 3.8 & 6.5 & 11.2 \\
Granite 3.1 8B & 24.5 & 33.0 & 26.2 & 24.5 & 34.0 & 47.5 \\
\bottomrule
\end{tabular}
\caption{The same questions with no page (\%). Family-only questions are
harder before any page is shown, which is why the main measure compares each
question with itself.}
\label{tab:closed}
\end{table}

\section{Examples of Evidence-Based Picks}

\begin{table}[h]
\centering
\small
\setlength{\tabcolsep}{3pt}
\begin{tabular}{>{\raggedright\arraybackslash}p{1.85cm}>{\raggedright\arraybackslash}p{2.35cm}>{\raggedright\arraybackslash}p{2.2cm}c}
\toprule
Drug & Graph answer & Evidence pick & Right? \\
\midrule
Lovastatin & familial hyperlipidemia & familial hyperlipidemia & yes \\
Tetrabenazine & juvenile Huntington disease & juvenile Huntington disease & yes \\
Moxifloxacin & Klebsiella pneumonia & Legionnaires' disease & no \\
Carboplatin & malignant germ cell tumor of ovary & ovarian carcinoma & no \\
\bottomrule
\end{tabular}
\caption{Examples of evidence-based picks within the given family (indication
edges). Misses often choose a better-studied member of the family.}
\label{tab:examples}
\end{table}

\end{document}
